\begin{longtable}[c]{|l|m{10cm}|}
    \caption{Real Number FFT Instructions}
    \label{tab:pie-real-fft-instructions}\\
    \hline
    \rowcolor{lightgray}
    \textbf{Instructions} & \textbf{Description} \\\hline
    \endfirsthead
    \rowcolor{lightgray}
    \textbf{Instructions} & \textbf{Description} \\\hline
    \endhead
    \endfoot
    
    \hyperref[instruction:EEFFTAMSS16LDINCPUAUP]{EE.FFT.AMS.S16.LD.INCP.UAUP} & Perform complex calculations and read 16-byte data from memory at the same time, then output 16-byte aligned data.
    \\\hline
    \hyperref[instruction:EEFFTAMSS16LDINCP]{EE.FFT.AMS.S16.LD.INCP} & Perform complex calculations and read 16-byte data from memory at the same time. Add 16 to address register.
    \\\hline
    \hyperref[instruction:EEFFTAMSS16LDR32DECP]{EE.FFT.AMS.S16.LD.R32.DECP} & Perform complex calculations and read 16-byte data from memory at the same time. Reverse the word order of read data. Add 16 to address register.
    \\\hline
    \hyperref[instruction:EEFFTAMSS16STINCP]{EE.FFT.AMS.S16.ST.INCP} & Perform complex calculation and write 16-byte data (consists of the data in AR and partial data segments in QR) to memory at the same time.
    \\\hline
    \hyperref[instruction:EEFFTVSTR32DECP]{EE.FFT.VST.R32.DECP} & Splice the QR register in 2-byte unit, shift the result and write this 16-byte data to memory.
    \\\hline
\end{longtable}